% Resumo (TCC2; lidera com os dois objetivos + a dissociação Φ/P_LLM).
% Parágrafo único; palavras-chave em main.tex via \palavraschave{}.

Este trabalho persegue dois objetivos pré-registrados: (i) formalizar métricas quantitativas para detectar a distorção comportamental induzida por RLHF e para medir a fidelidade comportamental em simulações sociais baseadas em LLMs, e (ii) usar tais métricas para testar se a fundamentação empírica em microdados socioeconômicos reais se propaga, através de agentes LLM ajustados por instrução, em comportamento coletivo realista. LLMs ajustados por instrução, quando empregados como agentes em simulações econômicas multiagentes, exibem uma anomalia sistemática --- taxas de cooperação que excedem tanto o equilíbrio de Nash quanto o comportamento humano de laboratório --- que formalizamos como o \textbf{Viés Cooperativo de RLHF} e operacionalizamos como a distância de variação total entre a distribuição de ações observada e o prior uniforme de ações, B\_RLHF = TV(π, π\_uniform); a fidelidade comportamental é medida por uma \textbf{Métrica de Realismo Comportamental (BRM ∈ [0,1])} composta. Ambas as métricas integram o \textbf{Arcabouço de Fundamentação Comportamental}, uma plataforma formalmente especificada BGF = (A, E, G, P, Φ, T), fundamentada em microdados da Pesquisa Social Europeia (ESS) Rodada 11 por meio de um mapa de síntese populacional Φ: D\_ESS → Profile. O achado central é uma \textbf{dissociação (Φ/P\_LLM)}: a fundamentação empírica, eficaz na camada de política baseada em regras, \textbf{não} se propaga através do canal de decisão do LLM. O teste pré-registrado de 10 sementes em N=100 encontra os braços fundamentado e não fundamentado estatisticamente indistinguíveis em cooperação (0,461 vs 0,455, MWU p = 0,91), Gini e riqueza --- a hipótese de fundamentação pré-registrada \textbf{não é corroborada na escala primária}, e relatamos esse resultado negativo de forma aberta. As métricas são precisamente o que revela a dissociação: na camada que lê a fundamentação diretamente, uma política determinística reproduz um coeficiente de Gini na faixa europeia (0,325, IC 95\% BCa [0,324, 0,325]; N=500, T=30, 10 sementes) e o gradiente cooperativo intercultural (ρ de Spearman = +1,000; World Values Survey r = +0,977; referencial de bens públicos de Herrmann--Thöni--Gächter ρ = +0,886, p = 0,033), estabelecendo que Φ é empiricamente válido enquanto a camada de decisão do LLM permanece ancorada no RLHF. Uma cascata cooperativa de RLHF dependente de escala (B\_RLHF amplificado 2,6--3,2× em N=500) e uma hipótese de profundidade de memória falsificada corroboram essa leitura de ancoragem. O arcabouço, as duas métricas e a dissociação são disponibilizados como artefato de código aberto (1.578 testes automatizados, reprodução por um único comando, testemunha criptográfica de reprodutibilidade), de modo que o resultado possa ser re-testado em qualquer cenário de LLM-como-agente sobre qualquer população empírica.
